\documentclass[border=4pt]{standalone}
\usepackage{amsmath,amssymb,mathtools,xcolor,varwidth}
\definecolor{simblue}{HTML}{1E6FE0}
\definecolor{simgreen}{HTML}{00A35A}
\begin{document}
\begin{varwidth}{28em}
\large\bfseries
Here stand two \emph{similar} figures: \textcolor{simblue}{the first figure $ABC$} and \textcolor{simgreen}{the second figure $DEF$}. To say that two figures are \emph{similar} means that they have exactly the same shape, though not necessarily the same size --- one may be a scaled copy of the other. All their corresponding angles are equal ($\angle A = \angle D$, $\angle B = \angle E$, $\angle C = \angle F$), and their corresponding sides run in the same proportion throughout. The matching outline shared by \textcolor{simblue}{the larger figure} and \textcolor{simgreen}{the smaller figure} --- their common shape --- is the heart of this lemma. Whatever we now prove about one such pair holds for \emph{all} similar figures, whether bounded by straight lines or by curves.
\end{varwidth}
\end{document}
